\section{The Controlled-Horizon Protocol}
\label{sec:protocol}

\paragraph{States.} We generate fights from fixed random seeds against five boss and elite encounters. Each fight advances by a short scripted prefix and is then frozen as a \emph{state}. Every state uses its own seed, so states are independent draws. A model never saw a state before the analysis plan was frozen. From the audited pool we release 1,760 states, 880 per character. Half of each character's states are \emph{sensitive}: the best move at $H{=}1$ and the best move at $H{=}8$ do not overlap. The other half are \emph{controls}, where some move is best at both horizons. The two sets are similar in encounter mix, number of legal moves, hand size, energy and health (Appendix~\ref{app:props}).

\paragraph{Prompt.} The model sees the full state as JSON: its health, energy, hand, the enemy's health and intent, every pile in its exact order and the remaining combat counters. It is told to choose ``the next action that maximizes the stated utility after exactly $H$ decision transitions.'' Utility is damage dealt minus health lost, plus 1,000 for a win and minus 1,000 for a loss. A decision is one card play or the end of a turn, so $H{=}8$ spans about two turns. We therefore test multi-turn lookahead within a fight. Payoffs that arrive after decision $H$ count for nothing, by design. A model that plays for the whole fight can therefore lose credit on some states. This is why the prompt states the horizon, the utility and the counting rule exactly, and a variant that defines ``decision transition'' explicitly changes nothing (Appendix~\ref{app:amendments}). It answers with one action in JSON. The prompts for different $H$ differ in that one number and nothing else, which we verify byte for byte.

\paragraph{Oracle.} For each state and horizon, depth-limited exhaustive search over the simulator returns the exact value $v_H(a)$ of every legal first action $a$. This value assumes best play for the remaining $H-1$ decisions. We score an answer by its normalized value
\[
\eqh{H}(a) = \frac{v_H(a) - \min_{a'} v_H(a')}{\max_{a'} v_H(a') - \min_{a'} v_H(a')},
\]
which is 1 for a best move and 0 for the worst. Illegal, unparseable or truncated answers score 0.

\paragraph{The main measure.} Let $a_1$ and $a_8$ be a model's answers to the same state at $H{=}1$ and $H{=}8$. The model's gain from the instruction is
\[
g = \eqh{8}(a_8) - \eqh{8}(a_1),
\]
both answers judged by the eight-decision standard. A model that ignores $H$ gives the same answer twice and earns $g{=}0$ on every state. Real gains come from changing the answer for the better. Our primary statistic is the difference between the mean $g$ on sensitive states and on control states. It is positive only if the model gains more where the horizon matters than where it does not. We test it per character with a 10,000-sample bootstrap over states at $\alpha = 0.025$, which keeps the two-character family at 0.05.

\paragraph{Lookahead use.} We also count how often the $H{=}8$ answer is a best eight-decision move on sensitive states. We compare that rate with an \emph{H-blind baseline}: the same rate for the model's own $H{=}1$ answer. A model that ignores $H$ matches its baseline exactly, and McNemar's exact test \citep{mcnemar1947} compares the two. We also report the \emph{greedy rate}, the share of answers that are best only at $H{=}1$.

\paragraph{Validity gates.} A model's run counts as confirmatory only if two conditions hold. First, at most 1\% of its answers were cut off by the 16,000-token output limit. Second, at most 2\% of states lost an answer to an execution failure. Runs that fail a gate are reported but do not enter any confirmatory test.

\paragraph{Pre-registration.} The states, prompts, scoring code, gates and analysis plan were frozen and fingerprinted before any model was run. Adding a model, a condition or a serving flag required a versioned amendment that may only add, never change. Appendix~\ref{app:amendments} lists every amendment. Analyses not in the frozen plan are labeled post hoc.
